% Requires \usepackage{amsmath,amssymb} in the document preamble.
% Methodology for the current arrival_final implementation.

\section{Methodology}
\label{sec:methodology}

\subsection{Hierarchical Architecture and Coordination Objective}

We consider $n$ car-like robots with fixed, individually assigned goals
$g_i\in\mathbb{R}^2$ and a common arrival deadline $T$. The objective is to
coordinate their motion so that all robots reach and remain at their goals
near $T$, with a small spread in arrival times and without inter-vehicle
collisions. The deadline is supplied at episode initialization using a nominal
travel-time estimate and randomized slack. Its construction uses the initial
missions of all robots; subsequent control decisions are decentralized given
this shared timing information.

The proposed hierarchy consists of a shared learned policy, an analytical
reference generator, and one nonlinear model predictive controller (NMPC) per
robot. The policy selects a local waypoint and preferred speed every
$\Delta_\pi=0.20\,\mathrm{s}$. Pure pursuit converts this action into a timed
reference by rolling out a kinematic bicycle model. NMPC updates every
$h=0.05\,\mathrm{s}$, using the latest measured state to track the reference
subject to input bounds and predicted separation constraints. The training
plant advances at $0.01\,\mathrm{s}$ intervals.

Let $\tau_i$ denote the recorded arrival time of robot $i$. Synchronized
completion requires all robots to satisfy the arrival condition, no collision,
and
\begin{equation}
    \max_i|\tau_i-T|\leq\epsilon_T,
    \qquad
    \max_i\tau_i-\min_i\tau_i\leq\epsilon_S,
\end{equation}
where $\epsilon_T=\epsilon_S=0.60\,\mathrm{s}$. These conditions distinguish
proximity to the deadline from synchronization between robots.

\subsection{Local Observations and Neighbour Selection}

The state of robot $i$ is
\begin{equation}
    x_i=[p_{x,i},p_{y,i},\psi_i,v_i,\delta_i]^\top,
\end{equation}
where $p_i=[p_{x,i},p_{y,i}]^\top$ is the rear-axle position, $\psi_i$ is
heading, $v_i$ is forward speed, and $\delta_i$ is steering angle. Let
$R_i=R(\psi_i)$ be the rotation from the robot frame to the world frame.
The goal is represented locally as $\widetilde g_i=R_i^\top(g_i-p_i)$.

Each robot receives a ten-dimensional ego observation containing its local
goal coordinates, speed, steering angle, remaining time $T-t$, deadline $T$,
initial goal distance $d_i^0$, arrival flag, required average speed, and
schedule error. With $d_i(t)=\|g_i-p_i(t)\|$, define
\begin{equation}
    q_i(t)=1-\frac{d_i(t)}{d_i^0},
    \qquad \eta(t)=\min\!\left(\frac{t}{T},1\right).
\end{equation}
The required-speed feature is
$d_i(t)/\max(T-t,\Delta_\pi)$, and the schedule feature is
$q_i(t)-\eta(t)$. Distances and times are scaled by $10\,\mathrm{m}$ and
$20\,\mathrm{s}$, respectively; speed and steering are scaled by their
limits. The normalized required speed is clipped to $[0,4]$, and the schedule
feature to $[-2,2]$.

The actor considers at most $K=4$ nearest neighbours within a sensing radius
$r_s=5\,\mathrm{m}$. Each selected neighbour $j$ is represented by
\begin{equation}
    o_{ij}^{\mathrm{nbr}}=
    \begin{bmatrix}
        R_i^\top(p_j-p_i)/r_s\\
        R_i^\top(\nu_j-\nu_i)/(2v_{\max})\\
        \sin(\psi_j-\psi_i)\\
        \cos(\psi_j-\psi_i)
    \end{bmatrix}\in\mathbb{R}^{6},
\end{equation}
where $\nu_i$ is the measured planar velocity. Missing neighbour slots are
padded and masked. The actor receives neither neighbour goals nor their
planned trajectories. The nominal sensing model supplies exact local states
within range and does not model occlusion or measurement noise.

\subsection{Attention-Based Shared Policy}

Separate two-layer multilayer perceptrons with $\tanh$ activations encode the
ego observation and each neighbour observation into $128$-dimensional vectors:
\begin{equation}
    e_i=f_{\mathrm{ego}}(o_i^{\mathrm{ego}}),
    \qquad h_{ij}=f_{\mathrm{nbr}}(o_{ij}^{\mathrm{nbr}}).
\end{equation}
The neighbour encoder shares parameters across all neighbour slots and robots.
An ego-conditioned attention mechanism computes
\begin{equation}
    q_i^{\mathrm{att}}=W_qe_i,
    \qquad k_{ij}=W_kh_{ij},
    \qquad \ell_{ij}=\frac{(q_i^{\mathrm{att}})^\top k_{ij}}{\sqrt{128}}.
\end{equation}
For the selected neighbour set $\mathcal{N}_i^K$, the weights and context are
\begin{equation}
    \alpha_{ij}=
    \frac{\exp(\ell_{ij})}{\sum_{m\in\mathcal{N}_i^K}\exp(\ell_{im})},
    \qquad c_i=\sum_{j\in\mathcal{N}_i^K}\alpha_{ij}h_{ij}.
\end{equation}
Invalid slots receive zero weight, and $c_i=0$ when no neighbours are present.
The neighbour embeddings themselves serve as the values in the weighted sum.

The concatenated representation $[e_i^\top,c_i^\top]^\top$ is passed through
an action head to produce the mean $\mu_i\in\mathbb{R}^3$ of a diagonal
Gaussian policy. During training,
\begin{equation}
    z_i\sim\mathcal{N}(\mu_i,\operatorname{diag}(\sigma^2)),
    \qquad a_i^{\mathrm{RL}}=\tanh(z_i),
\end{equation}
where the three log standard deviations are learned. Policy log-probabilities
include the Jacobian correction for the $\tanh$ transformation. Deterministic
execution uses $a_i^{\mathrm{RL}}=\tanh(\mu_i)$.

Shared encoding and attention pooling make the representation invariant to
permutations of the selected neighbours and independent of their count.
Nearest-neighbour selection bounds the actor's input workload; attention
learns how to weight the selected interactions. The same actor parameters
can therefore be used with different team sizes, although successful transfer
to larger teams remains an empirical question.

\subsection{Action Decoding and Timed Reference Generation}

The three bounded action components specify forward waypoint distance,
lateral offset, and preferred speed:
\begin{align}
    d_f &= 0.25+\frac{a_{i,1}^{\mathrm{RL}}+1}{2}(2.0-0.25),\\
    d_l &= 1.2a_{i,2}^{\mathrm{RL}},\\
    v_{\mathrm{pref}} &= \frac{a_{i,3}^{\mathrm{RL}}+1}{2}v_{\max}.
\end{align}
The world-frame waypoint is
\begin{equation}
    w_i=p_i+R_i[d_f,d_l]^\top.
\end{equation}
It remains fixed during one reference rollout. The third action is a speed,
not an arrival timestamp; temporal coordination is learned from the timing
observations and reward.

Starting from the measured state, the reference generator advances the
bicycle model for $N_r=24$ steps of duration $h$, producing $25$ states and
$24$ nominal controls over $1.2\,\mathrm{s}$. At each predicted state,
let $d_g$ and $d_w$ be the distances to the goal and waypoint. Braking-based
speed caps are
\begin{equation}
    v_g=\sqrt{2b_{\max}\max(d_g-\epsilon_g/2,0)},
    \qquad v_w=\sqrt{2b_{\max}d_w},
\end{equation}
where $b_{\max}$ is the maximum braking magnitude and
$\epsilon_g=0.25\,\mathrm{m}$ is the goal tolerance.

The pursuit target switches to the goal if the goal is already within
$\epsilon_g$, or if it lies ahead of the predicted vehicle and either
$d_g\leq d_w$ or $d_w<0.30\,\mathrm{m}$. This decision is reconsidered at
each rollout step. The target speed is $\min(v_{\mathrm{pref}},v_g)$ when
pursuing the goal and $\min(v_{\mathrm{pref}},v_g,v_w)$ otherwise.

For a pursuit target with local lateral coordinate $y_{\mathrm{loc}}$ and
squared distance $d^2$, the nominal steering and control commands are
\begin{align}
    \delta_{\mathrm{tar}}
    &=\operatorname{clip}\!\left(
        \operatorname{atan2}(2Ly_{\mathrm{loc}},\max(d^2,10^{-5})),
        -\delta_{\max},\delta_{\max}\right),\\
    u_k^{\mathrm{ref}}
    &=\operatorname{clip}_{\mathcal U}\!\left(
        \begin{bmatrix}
            4(v_{\mathrm{tar}}-v_k^{\mathrm{ref}})\\
            (\delta_{\mathrm{tar}}-\delta_k^{\mathrm{ref}})/h
        \end{bmatrix}\right).
\end{align}
The resulting references satisfy
$x_{k+1}^{\mathrm{ref}}=f_h(x_k^{\mathrm{ref}},u_k^{\mathrm{ref}})$.
They describe nominal motion at specific future times; the finite rollout
does not require the robot to reach its goal within $1.2\,\mathrm{s}$.

\subsection{Vehicle Prediction Model}

Reference generation and NMPC use a no-slip kinematic bicycle model with
wheelbase $L=0.35\,\mathrm{m}$. Its controls are longitudinal acceleration
and steering rate, $u=[a,\omega]^\top$, with bounds
\begin{equation}
    \mathcal U=
    \left\{[a,\omega]^\top:
        -2.0\leq a\leq1.5,\quad |\omega|\leq\pi/2\right\}.
\end{equation}
Acceleration is measured in $\mathrm{m/s^2}$ and steering rate in
$\mathrm{rad/s}$. Speed and steering are restricted to
$0\leq v\leq2.0\,\mathrm{m/s}$ and $|\delta|\leq27^\circ$.

The discrete transition saturates the speed and steering updates and uses
midpoint quantities for pose integration:
\begin{align}
    v^+&=\operatorname{clip}(v+ha,0,v_{\max}),\\
    \delta^+&=\operatorname{clip}(\delta+h\omega,-\delta_{\max},\delta_{\max}),\\
    \bar v&=(v+v^+)/2,\qquad \bar\delta=(\delta+\delta^+)/2,\\
    \Delta\psi&=h\bar v\tan(\bar\delta)/L,\\
    p^+&=p+h\bar v
       \begin{bmatrix}\cos(\psi+\Delta\psi/2)\\\sin(\psi+\Delta\psi/2)\end{bmatrix},\\
    \psi^+&=\psi+\Delta\psi.
\end{align}
Heading comparisons are wrapped modulo $2\pi$. This defines $f_h$; the
training plant uses the same integration rule at its smaller physics step.
Neither tyre-slip states nor actuator-lag states are included in the predictor.

\subsection{Constrained Nonlinear Model Predictive Control}

Each robot independently optimizes $N=16$ future controls over
$Nh=0.8\,\mathrm{s}$. A single-shooting formulation uses
$U=[u_0,\ldots,u_{N-1}]$ as the $32$ scalar decision variables and generates
predicted states recursively from the latest measurement. At control offset
$m\in\{0,1,2,3\}$ within a policy interval, the reference window is
$[x_m^{\mathrm{ref}},\ldots,x_{m+N}^{\mathrm{ref}}]$.

Let $r_k$ denote the corresponding state reference in this window. Define
$e_{p,k}=p_k-r_{p,k}$, $e_{v,k}=v_k-r_{v,k}$,
$e_{\delta,k}=\delta_k-r_{\delta,k}$, and
\begin{equation}
    e_{\psi,k}=\operatorname{atan2}\!\left(
       \sin(\psi_k-r_{\psi,k}),\cos(\psi_k-r_{\psi,k})\right).
\end{equation}
The objective is
\begin{equation}
\begin{aligned}
    J(U)=\sum_{k=0}^{N-1}\bigl(&
       8\|e_{p,k+1}\|^2+3e_{v,k+1}^2
       +0.5e_{\psi,k+1}^2+0.1e_{\delta,k+1}^2\\
       &+0.015\|u_k\|^2+0.04\|u_k-u_{k-1}\|^2\bigr)
       +16\|e_{p,N}\|^2+6e_{v,N}^2,
\end{aligned}
\end{equation}
where $u_{-1}$ is the previously applied control. The reference controls
initialize the optimizer; the objective penalizes control effort rather than
deviation from those reference controls. Consequently, optimized motion can
differ from pure pursuit even when the prediction models coincide.

For separation constraints, the vehicle footprint is covered by two
body-aligned discs. For body length $\ell_b$, width $w_b$, and rear-axle-to-centre
offset $\ell_c$, their common radius and longitudinal offsets are
\begin{equation}
    r_b=\sqrt{(\ell_b/4)^2+(w_b/2)^2},
    \qquad \ell_{1,2}=\ell_c\mp\ell_b/4.
\end{equation}
The configured dimensions are $\ell_b=0.55\,\mathrm{m}$,
$w_b=0.34\,\mathrm{m}$, and $\ell_c=0.16113\,\mathrm{m}$.
Disc centres are $c_i^a=p_i+\ell_a[\cos\psi_i,\sin\psi_i]^\top$.
Each sensed neighbour is extrapolated with constant speed and heading:
\begin{equation}
    \widehat p_{j,k}=p_j+khv_j[\cos\psi_j,\sin\psi_j]^\top,
    \qquad \widehat\psi_{j,k}=\psi_j.
\end{equation}
At every prediction node, all ego--neighbour disc pairs must satisfy
\begin{equation}
    \|c_{i,k}^{a}-\widehat c_{j,k}^{b}\|^2
       \geq (2r_b+0.08)^2,
    \qquad a,b\in\{1,2\},\quad k=1,\ldots,N.
\end{equation}
NMPC includes every neighbour within sensing range, independently of the
actor's nearest-neighbour truncation. The local problem minimizes $J(U)$
subject to these separation constraints, $u_k\in\mathcal U$, and the
prediction dynamics.

CasADi constructs the objective, constraint expressions, and their algorithmic
derivatives. Sequential least-squares programming (SLSQP) solves the local
nonlinear problem using a shifted previous solution as a warm start. The
reference controls provide an alternative initialization when the previous
candidate fails validation or warrants a lower-cost retry. A braking
initialization is also considered when the initial candidate is infeasible.
Only solver outputs passing finite-value, input-bound, and separation-residual
checks are accepted. Otherwise, bounded braking and steering straightening
are applied and the warm start is reset.

Only the first optimized control is applied. The robot is measured again
after $h$ seconds, and optimization repeats. The separation conditions constrain
sampled nominal predictions; constant-velocity forecasts and the fallback do
not establish continuous-time collision avoidance or recursive feasibility.

\subsection{Cooperative Reward and Arrival Detection}

Let $t^-$ and $t^+$ delimit a policy transition, with duration
$\Delta t=t^+-t^-$. Every robot receives the same cooperative reward:
\begin{equation}
\begin{aligned}
    r={}&\frac{3}{n}\sum_i\frac{d_i(t^-)-d_i(t^+)}{d_i^0}
       -\frac{0.5}{n}\sum_i|q_i(t^+)-\eta(t^+)|
       -0.01\Delta t\\
       &+r_{\mathrm{arrival}}+r_{\mathrm{end}}.
\end{aligned}
\end{equation}
Progress encourages goal-directed motion, while the schedule term encourages
distance-normalized progress to follow the common clock. This allows robots
with different trip lengths to learn different speed profiles.

Arrival requires the measured planar speed $s_i=\|\nu_i\|$ to satisfy
\begin{equation}
    s_i\leq0.08\,\mathrm{m/s},
    \qquad d_i+\frac{s_i^2}{2b_{\max}}\leq\epsilon_g
\end{equation}
continuously for a $0.20\,\mathrm{s}$ dwell interval. The recorded arrival
time is the beginning of that interval. Parked robots remain physical
obstacles, and arrival is revoked if the condition is subsequently violated.

For the set $\mathcal A_{\mathrm{new}}$ of robots receiving their first
arrival bonus in the transition,
\begin{equation}
    r_{\mathrm{arrival}}
    =\frac{5}{n}\sum_{i\in\mathcal A_{\mathrm{new}}}
       \exp\!\left[-\frac12\left(\frac{\tau_i-T}{\epsilon_T}\right)^2\right].
\end{equation}
The terminal reward is $-25$ for collision. For collision-free team completion,
it is
\begin{equation}
    r_{\mathrm{end}}=10\exp\!\left(
       -\frac{\max_i|\tau_i-T|}{\epsilon_T}
       -\frac{\max_i\tau_i-\min_i\tau_i}{\epsilon_S}\right).
\end{equation}
Otherwise, an episode reaching $T+3\,\mathrm{s}$ receives a penalty of
$-10$ times the fraction of robots that have not arrived. The terminal reward
is zero before termination. Collision, team completion, and timeout terminate
the episode. The same reward is used with both low-level controller variants.

\subsection{Training and Decentralized Execution}

The shared actor is trained using MAPPO with centralized training and
decentralized execution. A centralized critic encodes each robot's global
mission features, mean-pools the resulting embeddings into a team context,
and combines that context with each robot's embedding to predict one value
per robot. Its inputs include absolute states and goals, timing information,
and arrival and dwell history. These inputs are used only during training.

On-policy rollouts are collected through the complete control hierarchy,
including NMPC when training the full method. The critic supplies GAE value
estimates; the actor is updated using the PPO clipped surrogate, together
with value regression and entropy regularization. Actions selected by active
robots contribute to the policy loss, including actions that lead to arrival.
True task terminations suppress value bootstrapping, whereas artificial rollout
boundaries retain the next-state bootstrap and terminate the GAE trace.
Learning uses policy gradients through sampled transitions and does not
differentiate through pure pursuit, the NMPC solver, or the plant.

The training configuration uses three robots and a curriculum that broadens
from converging missions to mixed random, crossing, converging, and parallel
missions. Initial positions, goals, headings, speeds, and deadline slack are
randomized according to the active curriculum stage. The final protocol
excludes the dedicated extreme unequal-trip scenario family while retaining
ordinary variation in trip lengths. Scenario bounds define the sampling
region rather than physical walls. Training uses the kinematic bicycle plant.

At execution, the critic is discarded and every robot independently applies
the shared actor to its local observation. All controls for an interval use
the same pre-integration state snapshot. There is no leader, centralized
online planner, or exchange of neighbour goals and future plans. Controller
state and references are updated from measurements, while policy parameters
remain fixed.

For the pure-pursuit variant, the first four nominal reference controls are
applied over each policy interval. In the full method, these are replaced by
four NMPC solves using updated measurements. Adding NMPC after training a
pure-pursuit policy preserves the learned actor but changes the low-level
closed-loop dynamics; it is therefore distinct from training with NMPC inside
the environment. The same interface permits subsequent evaluation with a
different plant, such as MuJoCo, while retaining the nominal predictor and
fixed actor. Successful simulation transfer and larger-team performance must
be established through evaluation.
